\section{CNN 的构建模块}

本节课由 Zane Durante 主讲，分为两个主要部分：（1）如何构建 CNN——包括各类层组件和经典架构设计；（2）如何训练 CNN——包括权重初始化、学习率调度、数据增强等实用技巧。

\subsection{卷积层回顾}

\begin{figure}[H]
\centering
\includegraphics[width=0.95\textwidth]{slides-images/slide-001.jpg}
\caption{卷积层回顾：滤波器在输入图像上滑动，在每个位置计算内积产生激活图\protect\footnotemark}
\end{figure}
\footnotetext{Slides 第2页。}

卷积层是上节课的核心内容，这里简要回顾：
\begin{itemize}
    \item 每个滤波器的深度匹配输入通道数（如 RGB 图像的深度为 3）
    \item 滤波器在空间维度上滑动，在每个位置计算内积 + 偏置
    \item 每个滤波器产生一个二维激活图
    \item $C_{out}$ 个滤波器产生深度为 $C_{out}$ 的输出张量
    \item 激活图后通常接 ReLU 等非线性激活函数
\end{itemize}

\subsection{池化层回顾}

\begin{figure}[H]
\centering
\includegraphics[width=0.85\textwidth]{slides-images/slide-002.jpg}
\caption{池化层：$2 \times 2$ 最大池化（步幅 2）对空间维度进行下采样\protect\footnotemark}
\end{figure}
\footnotetext{Slides 第3页。}

池化层对每个通道独立地进行空间下采样。最常用的是 $2 \times 2$ 最大池化（stride 2），也有平均池化。

\subsection{CNN 构成总览}

\begin{figure}[H]
\centering
\includegraphics[width=0.95\textwidth]{slides-images/slide-003.jpg}
\caption{CNN 的完整组件：卷积层、池化层、全连接层（已学），归一化层、Dropout、激活函数（本节课）\protect\footnotemark}
\end{figure}
\footnotetext{Slides 第4页。}

到目前为止，课程已经介绍了三种基本层：
\begin{itemize}
    \item \textbf{卷积层}（Convolution）—— 提取空间特征
    \item \textbf{池化层}（Pooling）—— 空间下采样
    \item \textbf{全连接层}（Fully Connected）—— 矩阵乘法 + 偏置
\end{itemize}

本节课将补充三种重要组件：
\begin{itemize}
    \item \textbf{归一化层}（Normalization）—— 稳定训练
    \item \textbf{Dropout} —— 正则化
    \item \textbf{激活函数}（Activation Functions）—— 非线性
\end{itemize}

\subsection{本章小结}

CNN 由六类核心组件构成：卷积层负责提取特征，池化层负责下采样，全连接层负责最终映射，归一化层稳定训练过程，Dropout 提供正则化，激活函数引入非线性。理解每个组件的作用和组合方式是设计高性能 CNN 的基础。

%% ========== Part 2: 归一化层 ==========

